\section{Related Work}
\label{sec:related_work}

\para{Overrefusal and safety--helpfulness tradeoffs.}
Safety alignment via RLHF \citep{ouyang2022training} and constitutional AI \citep{bai2022constitutional} is known to produce overrefusal --- the rejection of safe prompts that superficially resemble harmful ones.
\citet{cui2024orbench} construct \orbench, an 80K-prompt benchmark that reveals a Spearman correlation of 0.89 between safety and overrefusal rates across 32 models, establishing overrefusal as a near-universal structural tradeoff.
\xstest \citep{rottger2024xstest} uses contrast pairs (safe and unsafe variants of the same topic) to identify exaggerated safety behaviors, while \sorrybench \citep{xie2024sorrybench} provides fine-grained evaluation across 44 safety categories.
These benchmarks measure \emph{explicit} refusal as a binary outcome.
Our work differs in two ways: we measure \emph{implicit} avoidance (hedging, disclaimers) as a continuous signal, and we target topics outside documented safety categories rather than borderline cases within them.

\para{Hidden refusal boundaries.}
\citet{rager2025discovering} introduce the Iterated Prefill Crawler (\ipc), which enumerates topics a model refuses to discuss without access to training data.
Applied to \llama and \deepseek, \ipc retrieves 31 of 36 known safety-trained topics and discovers additional hidden boundaries.
Notably, \deepseek exhibits ``thought suppression'' --- empty reasoning traces before refusing politically sensitive topics --- demonstrating internalized persona-level avoidance that generalizes beyond narrow training.
Our work complements \ipc by focusing on the \emph{non-refusal} manifestation of avoidance: models that engage with a topic but do so with excessive caution.

\para{Emergent misalignment and persona generalization.}
\citet{betley2025emergent} demonstrate that finetuning GPT-4o on 6{,}000 examples of covertly insecure code produces broadly misaligned behavior on unrelated tasks.
The mechanism is persona generalization: the training data implicitly encodes a ``deceptive assistant'' persona that the model generalizes beyond the training domain.
\citet{chen2025persona} provide a mechanistic account, identifying linear ``persona vectors'' in activation space that predict finetuning-induced behavioral shifts with $r = 0.76$--$0.97$.
These results establish that implicit behavioral patterns can generalize broadly from training data.
We hypothesize that the symmetric phenomenon occurs for caution: training data rich in hedging, disclaimers, and carefully qualified writing encodes a ``cautious assistant'' persona that generalizes to gray zone topics.

\para{Mechanistic interpretability of refusal.}
Recent work reveals that refusal is not a monolithic mechanism.
\citet{zhao2025harmfulness} show that harmfulness and refusal are encoded as separate linear directions in activation space, suggesting that avoidance is a separable behavioral layer overlaid on content understanding.
\citet{himelstein2025silenced} demonstrate that removing refusal via activation steering reveals latent biases that safety alignment masks rather than eliminates.
\citet{refusalnonlinear2025} find that refusal mechanisms are nonlinear and vary by architecture and layer depth.
These findings suggest that implicit avoidance, like explicit refusal, may be a distributed persona-level behavior rather than a simple decision boundary.

\para{Alignment costs.}
\citet{huang2025safetytax} document a ``safety tax'' on reasoning capabilities in large reasoning models, showing that safety alignment degrades performance on reasoning benchmarks.
\citet{perez2023sycophancy} establish that RLHF produces sycophancy --- another unintended behavioral pattern driven by training signal misalignment.
Together, these works show that alignment procedures produce systematic behavioral side effects beyond their intended safety function, consistent with our finding that implicit avoidance is a byproduct of training rather than a deliberate design choice.
